\section{Билет 70. Термодинамическое и статистическое определение энтропии. Энтропия и ее основные свойства. Теорема Нернста (третье начало термодинамики). Пример расчета изменения энтропии при смешивании газа при постоянной температуре.}

\begin{definition}
Существует два эквивалентных определения энтропии:
\begin{enumerate}
    \item \textbf{Термодинамическое (феноменологическое):} $dS = \frac{\delta Q_{\text{обр}}}{T}$ или $\Delta S = \int_1^2 \frac{\delta Q_{\text{обр}}}{T}$. Связывает энтропию с макроскопическим теплообменом и температурой.
    \item \textbf{Статистическое (микроскопическое):} $S = k \ln W$. Связывает энтропию с числом доступных микросостояний и степенью молекулярного хаоса.
\end{enumerate}
Оба подхода дают идентичные количественные результаты, но статистическое определение раскрывает физический смысл энтропии как меры неупорядоченности системы.
\end{definition}

\begin{property}[Основные свойства энтропии]
\begin{enumerate}
    \item \textbf{Функция состояния:} $\Delta S$ зависит только от начального и конечного состояний, а не от пути процесса.
    \item \textbf{Аддитивность (экстенсивность):} Энтропия составной системы равна сумме энтропий её частей.
    \item \textbf{Закон возрастания:} В изолированной системе $S$ не убывает: $\Delta S \ge 0$. Равновесию соответствует $S = S_{\max}$.
    \item \textbf{Мера необратимости:} $\Delta S > 0$ для необратимых процессов, $\Delta S = 0$ для обратимых.
    \item Не имеет абсолютного нуля в классической термодинамике (определяется с точностью до константы), но в статистической физике естественно нормируется через $W$.
\end{enumerate}
\end{property}

\begin{theorem}[Теорема Нернста --- третье начало термодинамики]
При стремлении температуры к абсолютному нулю энтропия любой термодинамической системы в равновесном состоянии стремится к постоянной величине, не зависящей от других параметров:
\begin{equation}
    \lim_{T \to 0} S(T) = S_0 = \text{const}. \label{eq:nernst_theorem}
\end{equation}
Планк предложил принять $S_0 = 0$ для идеально упорядоченных кристаллов. Тогда формула Больцмана даёт $W(T \to 0) = 1$, что соответствует единственному микросостоянию (полный порядок).
\par\noindent\textbf{Следствие (недостижимость абсолютного нуля):} Невозможно конечным числом термодинамических процессов охладить систему до $T = 0$ К.
\end{theorem}

\begin{example}[Расчёт $\Delta S$ при изотермическом смешивании газов]
Пусть в теплоизолированном сосуде разделённом перегородкой находятся два различных идеальных газа с количествами вещества $\nu_1$ и $\nu_2$ при одинаковой температуре $T$. Объёмы частей сосуда $V_1$ и $V_2$. После удаления перегородки газы смешиваются, занимая общий объём $V = V_1 + V_2$. Температура остаётся постоянной (внутренняя энергия идеального газа зависит только от $T$).
\par\noindent\textit{Решение:} Процесс смешивания необратим. Для расчёта $\Delta S$ мысленно заменим его обратимым изотермическим расширением каждого газа до объёма $V$.
\begin{equation}
    \Delta S_1 = \nu_1 R \ln\frac{V}{V_1}, \quad \Delta S_2 = \nu_2 R \ln\frac{V}{V_2}.
\end{equation}
Полное изменение энтропии системы:
\begin{equation}
    \Delta S_{\text{смеш}} = \Delta S_1 + \Delta S_2 = R \left( \nu_1 \ln\frac{V_1+V_2}{V_1} + \nu_2 \ln\frac{V_1+V_2}{V_2} \right). \label{eq:mixing_entropy}
\end{equation}
Так как $V > V_1$ и $V > V_2$, логарифмы положительны, следовательно $\Delta S_{\text{смеш}} > 0$. Это подтверждает необратимость процесса диффузии и рост энтропии при увеличении доступного объёма для каждой из подсистем.
\end{example}

\begin{remark}
Если смешиваются \textbf{одинаковые} газы, макроскопически состояние системы не меняется, и $\Delta S = 0$ (парадокс Гиббса). Разрешение парадокса даётся квантовой статистикой, учитывающей тождественность частиц: перестановка одинаковых молекул не создаёт нового микросостояния ($W$ не меняется, $S$ остаётся постоянной).
\end{remark}
